\documentclass[11pt]{article}
\usepackage[T1]{fontenc}
\usepackage{lmodern,microtype}
\usepackage{amsmath,amssymb,amsthm,mathtools}
\usepackage[letterpaper,margin=1in]{geometry}
\usepackage{booktabs,array,enumitem,longtable}
\usepackage{xcolor}
\usepackage[colorlinks=true,linkcolor=blue!55!black,citecolor=blue!55!black,urlcolor=blue!55!black]{hyperref}
\newcommand{\F}{F}
\newcommand{\K}{K}
\newcommand{\A}{\mathbb A}
\newcommand{\C}{\mathbb C}
\newcommand{\Z}{\mathbb Z}
\newcommand{\GL}{\operatorname{GL}}
\newcommand{\Irr}{\operatorname{Irr}}
\newcommand{\Nrd}{\operatorname{Nrd}}
\newcommand{\N}{\operatorname{N}}
\newcommand{\Hom}{\operatorname{Hom}}
\newcommand{\Rad}{\operatorname{Rad}}
\newcommand{\reg}{\mathrm{reg}}
\newcommand{\gen}{\mathrm{gen}}
\newcommand{\Ptag}{\textnormal{\textbf{[P]}}}
\newcommand{\Etag}{\textnormal{\textbf{[E]}}}
\newcommand{\Rtag}{\textnormal{\textbf{[R]}}}
\newcommand{\Otag}{\textnormal{\textbf{[O]}}}
\newcommand{\Ctag}{\textnormal{\textbf{[C]}}}
\newcommand{\tilG}{\widetilde G}
\newcommand{\tilH}{\widetilde H}
\newcommand{\tilT}{\widetilde T}
\newcommand{\wt}{\widetilde}
\theoremstyle{plain}
\newtheorem{theorem}{Theorem}[section]
\newtheorem{lemma}[theorem]{Lemma}
\newtheorem{proposition}[theorem]{Proposition}
\newtheorem{corollary}[theorem]{Corollary}
\theoremstyle{definition}
\newtheorem{definition}[theorem]{Definition}
\newtheorem{remark}[theorem]{Remark}
\title{Toward a metaplectic Jacquet--Langlands transfer, IV:\\
Wild Local Correspondence and Global Trace Formula Comparison}
\author{David Joyner\\[3pt]
\normalsize With mathematical and editorial assistance from ChatGPT (OpenAI)}
\date{November 7, 2026\quad---\quad Draft 1}
\begin{document}
\maketitle
\begin{abstract}
This continuation of Parts I--III studies the two unresolved steps in a
metaplectic Jacquet--Langlands transfer: integrality of the genuine
character transform at places dividing the covering degree, and
compatibility of the local transfer with a global trace-formula
comparison.  We formulate a precise integral-character-lattice
criterion for the wild local problem.  For the standard
Kazhdan--Patterson diagonal-torus cocycle, we compute the commutator
radical under the scalar-centrality hypothesis and obtain an explicit
good-element criterion for regular split diagonal elements, at both
tame and wild places.  We explain why this does not yet settle the
nonsplit-torus case, global geometric coefficients, or spectral
separation.  The latter questions are organized into verifiable proof
obligations.  No complete wild character correspondence or global
Jacquet--Langlands theorem is claimed.
\end{abstract}
\tableofcontents

\section{Introduction: a continuation with explicit proof boundaries}\label{sec:intro}
The classical Jacquet--Langlands correspondence relates representations
of $D^\times$ to essentially square-integrable representations of
$\GL_r(F)$, where $D/F$ is a central division algebra of degree $r$
\cite{DKV}.  A metaplectic variant involves genuine representations of
central extensions of these groups, and is not obtained merely by
pulling classical representations back to the covering groups.

The preceding parts of this series develop (I) the specific covering
cocycles, their centers, and good elements; (II) the local and global
matching of orbital integrals, under stated hypotheses; and (III) a
normalized character transform on the division-algebra side.  In
particular, Part III constructs an isometry
\begin{equation}\label{eq:LD}
 L_D:\mathcal C_\theta(D^\times)\longrightarrow
 \mathcal C_{\widetilde\theta}(\widetilde{D^\times})
\end{equation}
of invariant smooth functions with fixed central characters, proves a
tame character result, and constructs an adjoint orbital-integral
transfer, subject to the exact Part I cocycle and measure conventions.
At wild places the corresponding assertion on individual irreducible
characters remains open.

\paragraph{Notation.}
Locally $F$ is a nonarchimedean field of characteristic zero containing
$\mu_n$, $n\geq 1$ is odd, $r\geq 2$ is prime, and $(r,n)=1$.  Put
$H=D^\times$, $H^1=\ker\Nrd$, $Z=F^\times$ and
$\widetilde H=\widetilde{D^\times}$.  A genuine character of a cover
restricts to a prescribed faithful character of $\mu_n$.  We use the
central-character and star-map conventions of Part III.  When discussing
$G=\GL_r(F)$, $T$ denotes its diagonal torus and $Z_G$ its scalar
center.  Globally $K$ is a number field containing $\mu_n$.

\paragraph{Mathematical status labels.}
\textbf{[E]} means imported from a cited source or an earlier part of
the series; \textbf{[P]} is proved here under its explicit hypotheses;
\textbf{[R]} is a proved \emph{reduction}, not a solution of the original
problem; \textbf{[O]} is open; and \textbf{[C]} is a conditional
implication.  The exact Part I cocycle and the associated measure
normalizations have not been independently checked in this draft;
Claude's proposed cross-part review is therefore a particularly
valuable prerequisite before making claims stronger than those below.

\begin{center}\small
\begin{tabular}{@{}p{0.39\linewidth}p{0.48\linewidth}p{0.08\linewidth}@{}}
\toprule
Topic & Result in this draft & Status\\\midrule
Wild $D^\times$ character integrality & Character-lattice criterion and computable coefficients & [R]\\
Split diagonal-torus good elements & Exact radical under the specified torus cocycle & [P]\\
Good elements for nonsplit tori & Precise remaining local test & [O]\\
Global power classes modulo scalars & Almost-everywhere Hasse principle, recalled from III & [E]\\
Global trace-formula coefficients & Measure and orbit accounting checklist & [O]\\
Spectral separation and the local JL map & Conditional route, not a bijection & [C]\\
\bottomrule
\end{tabular}
\end{center}

\section{Wild-place character integrality as an integral-lattice problem}\label{sec:wild}
At residue characteristic $p\mid n$, Part III gives a unitary
transformation of invariant class functions but does not establish that
it maps every irreducible character to a genuine irreducible character.
Here is a precise way to isolate that missing assertion.

\subsection{The character matrix}\label{subsec:matrix}
Fix the central character $\theta$ of $H$ and the corresponding
genuine central character $\widetilde\theta$ of $\widetilde H$.
Let
\[
\mathcal R_\theta(H)=\bigoplus_{\rho\in\Irr_\theta(H)}\Z\Theta_\rho,
\qquad
\mathcal R_{\widetilde\theta}(\widetilde H)
=\bigoplus_{\tau\in\Irr_{\widetilde\theta}(\widetilde H)}
\Z\Theta_\tau
\]
be the \emph{integral virtual-character lattices}; the direct sums
mean finite integral combinations of characters.  For fixed unitary
central character, Schur orthogonality on the compact-mod-center
groups makes these irreducible characters orthonormal bases of the
respective invariant $L^2$ spaces.

Write
\begin{equation}\label{eq:coef}
 a_{\rho\tau}
 =\big\langle L_D\Theta_\rho,\Theta_\tau\big\rangle_{\widetilde H/\widetilde Z},
 \qquad L_D\Theta_\rho=\sum_\tau a_{\rho\tau}\Theta_\tau.
\end{equation}
The coefficients are \emph{complex} a priori.  With the regularity and
smoothness assertions of Part III, each expansion is finite, but this
finiteness should not be confused with integrality.

\begin{proposition}[\Rtag\ The integral-character-lattice criterion]\label{prop:lattice}
Assume the unitarity, smoothness, and finite character expansions of
Part III.  The following are equivalent:
\begin{enumerate}[label=(\roman*)]
 \item $L_D$ carries each irreducible character of $H$ to an
 irreducible genuine character of $\widetilde H$ and induces a bijection
 on their sets;
 \item $L_D\mathcal R_\theta(H)=
 \mathcal R_{\widetilde\theta}(\widetilde H)$;
 \item every coefficient $a_{\rho\tau}$ and every coefficient of
 $L_D^{-1}\Theta_\tau$ is an integer.
\end{enumerate}
In (iii) it is enough to prove integrality of the coefficients
$a_{\rho\tau}$ if the inverse matrix is the adjoint matrix and its
columns are the same finite families; then the inverse coefficients are
complex conjugates of integers and hence are integers.
The positivity at the identity specified by Part III determines the
positive sign in the resulting signed permutations.
\end{proposition}
\begin{proof}
(i) implies (ii) and (iii).  Condition (ii) says that $L_D$ and its
inverse preserve the indicated integral bases, and therefore implies
(iii).  Under (iii), the $\ell^2$-norm of every row of the unitary
matrix $A=(a_{\rho\tau})$ is one.  A finitely supported integral row
of norm one has exactly one nonzero entry, equal to $+1$ or $-1$.
The same argument for the inverse implies that these assignments
produce a bijection of the bases.  Finally, the identity-value formula
of Part III gives a strictly positive value for
$L_D\Theta_\rho$ at the identity, so the sign is positive.
\end{proof}

\begin{remark}[Why the inverse condition matters]
For an arbitrary map between free abelian groups, mapping a basis to
integral combinations need not imply surjectivity.  Here unitarity,
together with the basis expansion, is crucial: an integral unit vector
is already a signed basis vector.  The additional inverse statement
ensures that no target irreducible is omitted.
\end{remark}

\subsection{Numerical necessary conditions and a finite computation}\label{subsec:numerical}
Part III identifies the positive normalization
\begin{equation}\label{eq:cd}
 c_D=|n|_F^{-(r^2-1)/2},\qquad
 (L_D\Theta_\rho)(1)=c_D\dim\rho.
\end{equation}
Accordingly, if $L_D\Theta_\rho=\Theta_\tau$, then
\begin{equation}\label{eq:dim}
 \dim\tau=c_D\dim\rho.
\end{equation}
This condition is necessary, not sufficient.  For example, at
$F=\mathbb Q_3(\zeta_3)$, $r=2$, $n=3$, one has $|3|_F=3^{-2}$,
so $c_D=27$.  The dimension of a purported image must be divisible
by $27$.  The covered-torus argument in Part III supplies only a weaker
factor of $3$ for every genuine irreducible; it does not supply an
example that violates the stronger factor.

\begin{proposition}[\Rtag\ Finite-level test for a fixed representation]\label{prop:finite-test}
Fix $\rho\in\Irr_\theta(H)$ and assume the smoothness assertion for
$L_D\Theta_\rho$ from Part III.  Only finitely many
$a_{\rho\tau}$ can be nonzero, and integrality of the image character
is decided by the finite collection of equalities
\[
 a_{\rho\tau}\in\Z\qquad
 (a_{\rho\tau}\ne0).
\]
Each coefficient can be computed by finite character sums after
choosing sufficiently small compact open normal subgroups compatible
with the given central characters.
\end{proposition}
\begin{proof}
A smooth invariant function with fixed central character on a
compact-mod-center group is determined at a sufficiently small normal
level of the resulting compact projective quotient.  Its character expansion lies in the finite
character space of that finite twisted quotient.  Integration against
the normalized Haar measure becomes a finite sum on this quotient.
The coefficient formula is exactly \eqref{eq:coef}.
\end{proof}

\begin{remark}[A falsifiable research direction]\label{rem:counter}
A single $\rho$ for which one coefficient in \eqref{eq:coef} is
nonintegral disproves irreducibility preservation by the \emph{specified}
$L_D$.  It need not disprove the existence of another wild local
correspondence.  Conversely, checking finitely many representations
at finitely many levels does not prove the statement at every depth.
\end{remark}

\subsection{Why the tame Adams proof cannot be repeated}\label{subsec:adams}
For a finite group $Q$, the Adams operation satisfies
$\psi^n\chi(g)=\chi(g^n)$.  When $n$ is prime to $|Q|$,
cyclotomic Galois conjugation makes this a permutation of irreducible
characters; this is the ingredient used in the tame case of Part III.
It fails when $p\mid n$ and $Q$ has $p$-torsion.

\begin{remark}[An elementary obstruction]
For $Q=C_p$ and $n=p$, every $g\in Q$ satisfies $g^p=1$, hence
$\psi^p\chi(g)=\chi(1)$.  In particular all linear characters have the
same Adams image, the trivial character.  Thus the tame Adams
argument is not available at a wild place without new structure.
\end{remark}

\section{A split-place theorem: the commutator radical}\label{sec:split}
The preceding discussion concerns the division side.  A second
bottleneck is to characterize \emph{good} regular elements of the
metaplectic cover at places where the division algebra is split.
The first case, regular elements in the split diagonal torus, admits
an exact calculation.

\subsection{The standard torus cocycle and its matrix}\label{subsec:cocycle}
Let $G=\GL_r(F)$, with diagonal torus $T\simeq(F^\times)^r$ and
scalar center $Z_G\simeq F^\times$.  Assume the restriction of the
chosen cocycle to $T$ is represented by the standard
Kazhdan--Patterson convention
\begin{equation}\label{eq:sigma}
 \sigma_c(t,u)=(\det t,\det u)_n^c
 \prod_{1\le i<j\le r}(t_i,u_j)_n,
 \quad t=\operatorname{diag}(t_i),\quad u=\operatorname{diag}(u_i).
\end{equation}
Here $(\cdot,\cdot)_n$ is the $n$th Hilbert symbol on $F$.  This
normalization is compatible with the scalar-centrality congruence
used in Part I,
\begin{equation}\label{eq:scalar}
 r-1+2rc\equiv0\pmod n.
\end{equation}
Equivalent cocycles give the same commutator pairing, but before
using this calculation globally the convention \eqref{eq:sigma}
needs to be checked against the exact local section used in Part I.
The Kazhdan--Patterson and later block-compatible formulas
are discussed in \cite{FK86,Takeda,ZouTypes}.  The simple-type
results in \cite{ZouTypes} assume $(p,n)=1$ and therefore do not
resolve the present wild problem.

Since $n$ is odd, the Hilbert symbol is alternating:
$(x,y)_n=(y,x)_n^{-1}$.  Directly taking the commutator of the two
cocycle lifts gives
\begin{equation}\label{eq:B}
 B_c(t,u):=[\widetilde t,\widetilde u]
 =\prod_{i,j=1}^r(t_i,u_j)_n^{A_{ij}},
 \quad A=-I_r+(1+2c)J_r,
\end{equation}
where $J_r$ is the all-ones matrix.  In particular the diagonal
entries of $A$ are $2c$ and the off-diagonal entries are $1+2c$.
Its row and column sums equal $r-1+2rc$, which vanish modulo $n$
under \eqref{eq:scalar}; accordingly the scalar lifts commute with
$\widetilde T$.

\begin{lemma}[\Ptag\ Radical of the commutator matrix]\label{lem:rad}
Assume $n$ is odd and $(r,n)=1$, and let
$X=F^\times/F^{\times n}$.  The bilinear pairing
\[
 X^r\times X^r\longrightarrow\mu_n,
 \qquad (x,y)\longmapsto\prod_{i,j}(x_i,y_j)_n^{A_{ij}}
\]
with $A$ as in \eqref{eq:B} and subject to \eqref{eq:scalar}
has radical exactly the diagonal subgroup
$\Delta(X)=\{(a,\ldots,a):a\in X\}$.
\end{lemma}
\begin{proof}
The Hilbert pairing $X\times X\to\mu_n$ is perfect by local
Kummer theory and local duality.  Hence the radical of the displayed
pairing is precisely $\ker(A:X^r\to X^r)$.
The group $X$ is killed by $n$.  Because $(r,n)=1$, multiplication by
$r$ is an automorphism of $X$, and every $x\in X^r$ decomposes
uniquely as
\[
 x=(a,\ldots,a)+w,
 \qquad \sum_iw_i=0.
\]
On the diagonal summand, $A$ acts by
$r-1+2rc\equiv0\pmod n$.  On the sum-zero summand it acts as
$-I_r$, since $J_rw=0$.  Thus it is invertible on that summand and
its kernel is precisely $\Delta(X)$.
\end{proof}

\begin{theorem}[\Ptag\ Good regular elements of the split diagonal torus]\label{thm:split-good}
Let $t=\operatorname{diag}(t_1,\ldots,t_r)$ be regular in
$G=\GL_r(F)$.  Assume \eqref{eq:sigma} and \eqref{eq:scalar}, with
$n$ odd and $(r,n)=1$.  A lift $\widetilde t$ is good, in the sense
that it commutes with the full preimage of $G_t=T$, if and only if
\begin{equation}\label{eq:goodcrit}
 t\in Z_GT^n.
\end{equation}
Equivalently,
\begin{equation}\label{eq:ratios}
 t_i/t_j\in F^{\times n}\quad\text{for all }i,j.
\end{equation}
This criterion applies at both tame and wild places.
\end{theorem}
\begin{proof}
For regular $t$, its centralizer in $G$ is $T$.  By the definition
of good elements, $\widetilde t$ is good if and only if its commutator
with every lift of $u\in T$ is trivial.  In view of \eqref{eq:B}, this
is exactly membership in the radical of the torus pairing.  Lemma
\ref{lem:rad} identifies the radical on $X^r$ with the diagonal.
The inverse image of this subgroup in $T$ is $Z_GT^n$.
Condition \eqref{eq:ratios} is equivalent to that assertion, since
$t_i/t_1$ being an $n$th power allows one to factor $t$ into a scalar
and a coordinatewise $n$th power.
\end{proof}

\begin{remark}[The star-map description]
Under the scalar-centrality hypothesis, a regular good diagonal
lift is a central lift times the $n$th power of a lift of another
diagonal element.  Thus the result matches the star-map description
used on the division side in Part III; it does not require the
$n$th-power map to be an automorphism of $T$.
\end{remark}

\begin{remark}[Concrete example: $r=2$, $n=3$]
Take $F=\mathbb Q_3(\zeta_3)$ and $c\equiv2\pmod3$.  Then
$1+4c\equiv0\pmod3$ and
\[
 A=\begin{pmatrix}1&-1\\-1&1\end{pmatrix}\pmod3,
 \qquad
 B_c(t,u)=(t_1/t_2,u_1/u_2)_3.
\]
Perfection of the cubic Hilbert pairing shows that a regular
$\operatorname{diag}(t_1,t_2)$ is good exactly when
$t_1/t_2\in F^{\times3}$; the residue characteristic dividing $n$
does not alter this equivalence.
\end{remark}

\subsection{Exactly what the split calculation does not establish}
A \emph{split place for $D$} means $D\otimes_KK_v\simeq M_r(K_v)$;
it does not mean that every elliptic or regular maximal torus in
$\GL_r(K_v)$ is split.  The centralizer of a regular element can be
$\operatorname{Res}_{E_v/K_v}\mathbb G_m$ for a \emph{field} or a
product of fields $E_v$, not necessarily the diagonal torus.
Theorem~\ref{thm:split-good} therefore does not by itself identify
all good regular classes in the local split group.

\begin{definition}[\Otag\ Nonsplit-torus problem]\label{def:nonsplit}
For each degree-$r$ \emph{\'etale} $F$-algebra $E$ embedded as a
regular maximal torus $T_E\subset\GL_r$, compute the radical of the
commutator pairing on $p^{-1}(T_E(F))$, with the exact local cocycle of
Part I.  Determine whether it is the preimage of
$F^\times T_E(F)^n$, or whether additional local norm or cohomological
classes occur.  The answer must be uniform in $E$ and include the
places where $p\mid n$.
\end{definition}

A possible route is to express the commutator via restriction and
corestriction of Hilbert symbols on the factors of $E$.  At a split
place this requires a computation on each rational torus, not merely
a calculation after extending scalars to an algebraic closure.
A formula valid after scalar extension need not automatically imply
an equality of radicals on $F$-rational points.

\section{Global power classes, rational orbits, and their limits}\label{sec:global}
The split diagonal calculation suggests a connection between good
elements and $n$th powers modulo scalars.  Part III established a
useful global arithmetic result in precisely this direction.

\begin{proposition}[\Etag\ Almost-everywhere Hasse principle, Part III]\label{prop:hasse}
Let $E/K$ be a field extension of degree $r$, let $\mu_n\subset K$,
and assume $(r,n)=1$.  For any finite set $\Sigma$ of places of $K$,
the localization map
\begin{equation}\label{eq:hasse}
 \frac{E^\times}{K^\times E^{\times n}}
 \longrightarrow
 \prod_{v\notin\Sigma}
 \frac{(E\otimes_KK_v)^\times}
 {K_v^\times(E\otimes_KK_v)^{\times n}}
\end{equation}
is injective.
\end{proposition}
\begin{proof}[Proof recalled for completeness]
Suppose the class of $t\in E^\times$ is trivial in every local
factor outside $\Sigma$.  For such $v$ write $t=a_vy_v^n$ with
$a_v\in K_v^\times$ and $y_v\in(E\otimes_KK_v)^\times$.
Choose integers $u,k$ such that $ru+nk=1$.  Put
$a=\N_{E/K}(t)^u\in K^\times$.  Taking local norms and rearranging
gives
\[
 \frac ta=\left(a_v^k\frac{y_v}{\N(y_v)^u}\right)^n.
\]
Thus $t/a$ is an $n$th power in every completion of $E$ above
$v\notin\Sigma$.  If $t/a$ were not an $n$th power in $E$, the
Kummer extension $E((t/a)^{1/n})/E$ would be nontrivial because
$\mu_n\subset E$.  Chebotarev would then give infinitely many
unramified places of $E$ at which $t/a$ is not a local $n$th
power, a contradiction.  Hence $t\in K^\times E^{\times n}$.
\end{proof}

\begin{remark}[No automatic global cover comparison]
The proposition tests local powers modulo scalars at \emph{almost all
places}.  It cannot be invoked from the good-element condition unless
that condition has first been identified with the local power
condition for the \emph{actual} regular centralizer at those places.
Theorem~\ref{thm:split-good} proves one case, not the needed
uniform statement.  Global reciprocity for the covering cocycle and
compatibility of adelic star lifts are additional questions.
\end{remark}

\subsection{What Kottwitz's framework already resolves}
Let $D/K$ be a central simple algebra of degree $r$, and let a
regular elliptic element correspond to a degree-$r$ field $E/K$.
Its centralizer as an algebraic group is the induced torus
$T=\operatorname{Res}_{E/K}\mathbb G_m$.  Hilbert's Theorem 90 and
Shapiro's lemma give $H^1(K,T)=1$, so a represented stable regular
class does not split into multiple rational conjugacy classes.
Existence of an embedding $E\hookrightarrow D$ is equivalent to
$D\otimes_KE$ being split; the Brauer--Hasse--Noether theorem detects
this condition locally.

These are special simplifications, in the inner forms of $\GL_r$,
of the rational-versus-stable orbit and torus-transfer formalism of
Kottwitz \cite[\S\S7--9]{Kottwitz84}.  His formula for relative
Tamagawa numbers \cite[\S5]{Kottwitz84} supplies the correct context
for global volume factors.  None of these linear-group statements by
itself controls genuine conjugacy, the star map, or the adelic
covering cocycle.

\section{Geometric normalization: the next exact comparison}\label{sec:measures}
The trace formulas cannot be compared by multiplying the local
probability-measure constants from Part III without examining their
relation to the globally fixed Haar measures.

\subsection{Three distinct quantities}
At a division place $v$, Part III introduces
\[
 b_{D,v}=|n|_v^{-(r-1)/2}
 \quad\text{and}\quad
 c_{D,v}=|n|_v^{-(r^2-1)/2}.
\]
The former is the normalized elliptic class-function/orbital-integral
factor for the \emph{local probability quotient measure}; the latter
is its value-at-identity scaling.  Neither is automatically the
global geometric factor.  A third contribution is the conversion
between those local probability measures and the Haar measures used
in the adelic trace formula.  These three quantities must remain
separate in every local comparison.

The product formula $\prod_v|n|_v=1$ is true for $n\in K^\times$,
but does not by itself cancel the $b_{D,v}$: these constants occur
at the division places, whereas the complement includes split and
archimedean places with their own transfer factors.

\subsection{An explicit normalization ledger}\label{subsec:ledger}
For each rational regular elliptic class $\gamma$, the desired
geometric identity should be checked term by term with the following
data, using one globally fixed normalization:
\begin{center}\small
\begin{tabular}{@{}p{0.25\linewidth}p{0.63\linewidth}@{}}
\toprule
Input & Required verification\\\midrule
Local orbit & Is the lift good? Does the rational stable class have a matching lift?\\
Local orbital integral & Exact transfer factor, sign, power-map multiplicity, and Jacobian\\
Local quotient measure & Conversion of $T_v/Z_v$ probability measure to the adelic Haar measure\\
Global torus volume & $T(K)\backslash T(\A_K)$ factor with the center and connected components treated consistently\\
Global cocycle & Product/reciprocity law for the chosen lifts; dependence on $c_v$\\
Singular and identity terms & Correct treatment beyond regular elliptic classes, or a test-function hypothesis eliminating them\\
Spectral side & Absolute convergence or finite packet identity and a justified separation argument\\
\bottomrule
\end{tabular}
\end{center}

\begin{definition}[\Otag\ Geometric comparison problem]\label{def:geom}
Fix global measures and matched factorizable genuine test
functions.  Prove a trace-formula identity by identifying the
regular elliptic geometric terms, including their global centralizer
volumes and all local factors, and then handling or eliminating
singular terms.  The comparison is \emph{not} asserted by the
existence of local transfer operators alone.
\end{definition}

\begin{remark}[Available global trace formula]
Wen-Wei Li constructed an invariant trace formula for adelic
covers of connected reductive groups under specified hypotheses,
including relevant covers of $\GL_r$ \cite{Li14}.  This is an
important global tool, but it must be specialized to the present
cocycles and matched functions.  A theorem about the existence of a
trace formula is not the same as the required equality of the two
geometric sides.
\end{remark}

\section{Spectral separation and a proposed order of attack}\label{sec:spectral}
Even a geometric identity does not automatically produce a
bijection between local irreducible representations.  A global
spectral comparison will yield combinations of local characters,
possibly with multiplicities.  To recover an \emph{individual}
character, one needs a separation principle, a supply of suitable
globalizations, and control of every other local factor.

\begin{proposition}[\Rtag\ What a spectral identity must detect]\label{prop:spectral}
Let $\rho\in\Irr_\theta(D_v^\times)$ be fixed at a wild place.
Suppose a proved global identity, using test functions varied at $v$,
implies the exact equality of local invariant distributions
\[
 L_D\Theta_\rho=\sum_{\tau}m_\tau\Theta_\tau,
 \qquad m_\tau\in\Z,
\]
with only finitely many terms and the normalization of Part III.
Then $L_D\Theta_\rho$ equals a single irreducible genuine
character.  If the same conclusion holds for all $\rho$ and the
inverse coefficients are integral, the local wild correspondence is a
bijection.
\end{proposition}
\begin{proof}
This is Proposition~\ref{prop:lattice}, applied to one row of the
unitary character matrix.  The hypothesis that the spectral identity
supplies \emph{integer} coefficients is the essential additional
input: neither a geometric identity nor $L^2$ unitarity establishes
it without an argument.
\end{proof}

\paragraph{Preferred order of work.}
The next rigorous advance should be the nonsplit regular-torus
commutator calculation, followed by the unramified transfer and the
adelic normalization ledger.  In parallel, the finite-level
computation of \eqref{eq:coef} for the smallest wild example
$(r,n,F)=(2,3,\mathbb Q_3(\zeta_3))$ is a useful independent test:
one nonintegral coefficient would change the strategy immediately.
Only after these checks should the spectral comparison be promoted
to a claimed theorem.  Removal of the auxiliary place and global
non-cancellation hypotheses belongs after a correct base trace
identity is available.

\section{Glossary of specialized terms}\label{sec:glossary}
\begin{description}[style=nextline,leftmargin=2.5cm]
\item[Adams operation]
The operation $\psi^n$ on a representation ring, characterized on
characters of a finite group by $\psi^n\chi(g)=\chi(g^n)$.  It is
related to exterior powers in equivariant $K$-theory, and may produce
virtual rather than irreducible characters.
\item[Character lattice]
The free abelian group generated by irreducible characters with a
fixed central character; its elements are integral virtual characters.
\item[Good element]
A lift whose commutator with the full preimage of the ordinary
centralizer is trivial.  A genuine invariant distribution is
supported on the appropriate good locus.
\item[Hilbert symbol]
The local norm-residue pairing $(\cdot,\cdot)_n$ on
$F^\times/F^{\times n}$, valued in $\mu_n$.  It is perfect when
$\mu_n\subset F$.
\item[Commutator radical]
The subgroup of a torus on which the lifted commutator pairing is
trivial against every element of that torus.
\item[Split place / split torus]
A split place of $D$ is one where $D_v\simeq M_r(K_v)$; a split torus
is a torus isomorphic over $K_v$ to a product of copies of
$\mathbb G_m$.  These are different notions.
\item[Central character]
A prescribed scalar by which the center acts in an irreducible
representation.  Compatibility of these characters is required when
comparing representations of different covers.
\item[Integral virtual character]
A finite integer combination of irreducible characters, not
necessarily the character of an actual representation.
\item[Rational and stable conjugacy]
Conjugacy under $G(K)$ and under $G(\overline K)$, respectively.
Their difference is measured by Galois cohomology of centralizers.
\item[Induced torus]
For $E/K$ finite \'etale, $\operatorname{Res}_{E/K}\mathbb G_m$;
its $K$-points are $E^\times$ and its first Galois cohomology
vanishes by Shapiro and Hilbert 90.
\item[Kottwitz obstruction]
A cohomological obstruction to realizing locally transferred maximal
tori in a global inner form; in Kottwitz's setting it is expressed
using $H^2(K,T_{\mathrm{sc}})$.
\item[Tamagawa measure]
The canonical global measure built from compatible local algebraic
measures and convergence factors.  It need not agree with probability
Haar measure on a compact local quotient.
\item[Kummer class]
A class in $E^\times/E^{\times n}$ or its torus analogue;
Kummer theory converts it into a finite Galois-cohomology class.
\item[Spectral separation]
The deduction of individual local character identities from a
global equality of sums of automorphic representation traces.
\item[Local-to-global obstruction]
A class that disappears in every localization but remains nontrivial
globally.  The power-class kernel recalled here vanishes under the
stated hypotheses; covering-group compatibility may have additional
obstructions.
\end{description}

\section*{Acknowledgments and verification scope}
The author acknowledges mathematical discussions and editorial
assistance from ChatGPT (OpenAI).  The earlier parts benefited from
Claude (Anthropic), whom the author plans to consult separately for a cross-part
assessment of Parts I--III.  The new radical computation and its
application to split diagonal tori should be checked against the
exact cocycle conventions of Part I before publication.  No
unverified global metaplectic correspondence is claimed.

\begin{thebibliography}{99}
\bibitem{DKV} P. Deligne, D. Kazhdan and M.-F. Vign\'eras,
\emph{Repr\'esentations des alg\`ebres centrales simples $p$-adiques},
in \emph{Representations of Reductive Groups over a Local Field} (1984).
\bibitem{FK86} Y. Flicker and D. Kazhdan,
\emph{Metaplectic correspondence}, Publ. Math. IH\'ES \textbf{64}
(1986), 53--110.
\url{https://www.numdam.org/item/PMIHES_1986__64__53_0/}.
\bibitem{Takeda} S. Takeda, \emph{Remarks on metaplectic tensor products for covers of $\mathrm{GL}_r$},
Pacific J. Math. \textbf{290} (2017); see the scalar-center formula.
\url{https://msp.org/pjm/2017/290-1/pjm-v290-n1-p07-p.pdf}.
\bibitem{ZouTypes} J. Zou, \emph{Simple type theory for metaplectic
covers of $\mathrm{GL}(r)$ over a non-archimedean local field},
J. Inst. Math. Jussieu \textbf{24} (2025), no.~4, 1263--1335, \S4.
\url{https://doi.org/10.1017/S1474748025000015}.
\bibitem{Kottwitz84} R. E. Kottwitz, \emph{Stable trace formula:
cuspidal tempered terms}, Duke Math. J. \textbf{51} (1984),
611--650, \S\S5,7--9.
\url{https://doi.org/10.1215/S0012-7094-84-05129-9}.
\bibitem{Li14} W.-W. Li, \emph{La formule des traces pour les rev\^etements
de groupes r\'eductifs connexes. IV. Distributions invariantes},
Ann. Inst. Fourier \textbf{64} (2014), 2379--2448.
\url{https://doi.org/10.5802/aif.2915}.
\bibitem{Cohen} J. Cohen, \emph{Transfer of representations and
orbital integrals for inner forms of $\mathrm{GL}_n$},
Canad. J. Math. \textbf{70} (2018), 595--627.
\bibitem{Adams} J. F. Adams, \emph{Vector fields on spheres},
Ann. of Math. \textbf{75} (1962), 603--632.
\end{thebibliography}
\end{document}
